%=====================================================================
%  FAKERNI — End-of-Studies / Internship Report (PFE)
%  ESPRIT-style layout (modelled on a standard Tunisian PFE report).
%
%  HOW TO USE:
%   1. Upload the whole "report/" folder to Overleaf (this file + images/).
%   2. Set the compiler to pdfLaTeX.
%   3. Edit every \myXxx value in the "EDIT THESE" block below with your
%      real details (name, firm, supervisors, school, dates).
%   4. Drop your diagrams into images/ using the file names referenced by
%      each \figph{...} command; until then a labelled placeholder box
%      is shown so the document still compiles.
%=====================================================================
\documentclass[12pt,a4paper,oneside]{report}

\usepackage[utf8]{inputenc}
\usepackage[T1]{fontenc}
\usepackage[english]{babel}
\usepackage[a4paper,margin=2.5cm]{geometry}
\usepackage{graphicx}
\usepackage{float}
\usepackage{xcolor}
\usepackage{booktabs}
\usepackage{longtable}
\usepackage{array}
\usepackage{tabularx}
\usepackage{enumitem}
\usepackage{caption}
\usepackage{fancyhdr}
\usepackage{titlesec}
\usepackage{listings}
\usepackage[hidelinks]{hyperref}

\graphicspath{{images/}}

% ---------------------------------------------------------------------
% EDIT THESE  (all the personal / institutional details)
% ---------------------------------------------------------------------
\newcommand{\mySchool}{[Your School / University]}          % e.g. ESPRIT School of Business
\newcommand{\myReportType}{INTERNSHIP REPORT}               % or "END-OF-STUDIES PROJECT REPORT"
\newcommand{\myDegree}{Bachelor's Degree in [Your Programme]}
\newcommand{\mySpecialization}{[Your Specialization]}
\newcommand{\myProjectTitle}{Design and Development of ``Fakerni'': a Collaborative,
Real-Time and AI-Assisted Shared-List Platform for Groups}
\newcommand{\myAuthor}{Ghassen Chich}                        % <-- your full name
\newcommand{\myCompany}{[Host Company Name]}                 % <-- the firm
\newcommand{\myProSupervisor}{[Professional Supervisor]}     % supervisor at the firm
\newcommand{\myAcadSupervisor}{[Academic Supervisor]}        % supervisor at the school
\newcommand{\myPresentedOn}{[dd/mm/yyyy]}
\newcommand{\myPeriod}{[Month Year -- Month Year]}
\newcommand{\myAcademicYear}{2025 / 2026}
% ---------------------------------------------------------------------

% ---- Colours & code style ------------------------------------------
\definecolor{brand}{RGB}{13,110,120}
\definecolor{brandred}{RGB}{192,57,43}
\definecolor{codebg}{RGB}{245,247,248}
\definecolor{codekw}{RGB}{26,107,115}
\definecolor{codecm}{RGB}{120,124,126}

\lstset{
  backgroundcolor=\color{codebg},
  basicstyle=\ttfamily\footnotesize,
  keywordstyle=\color{codekw}\bfseries,
  commentstyle=\color{codecm}\itshape,
  breaklines=true,
  frame=single,
  rulecolor=\color{codebg},
  showstringspaces=false,
  columns=fullflexible,
}

% ---- Headers / footers ---------------------------------------------
\pagestyle{fancy}
\fancyhf{}
\fancyhead[L]{\small\itshape Fakerni --- \myReportType}
\fancyhead[R]{\small\thepage}
\renewcommand{\headrulewidth}{0.4pt}

% ---- Chapter title styling -----------------------------------------
\titleformat{\chapter}[display]
  {\normalfont\huge\bfseries\color{brand}}
  {\chaptertitlename\ \thechapter}{16pt}{\Huge}

% ---- Figure placeholder helper -------------------------------------
% \figph{filename.png}{Caption text}
% Includes images/filename.png if present, otherwise draws a labelled box.
\newcommand{\figph}[2]{%
  \begin{figure}[H]\centering
  \IfFileExists{images/#1}
    {\includegraphics[width=0.92\textwidth]{#1}}
    {\setlength{\fboxsep}{10pt}%
     \fbox{\parbox[c][4.2cm][c]{0.82\textwidth}{\centering
       \textit{[ Figure placeholder ]}\\[4pt]
       Add your image as \texttt{images/#1}\\[4pt]#2}}}
  \caption{#2}\end{figure}}

% Small logo helper (school / company)
\newcommand{\logoph}[2]{%
  \IfFileExists{images/#1}
    {\includegraphics[height=2.4cm]{#1}}
    {\setlength{\fboxsep}{6pt}\fbox{\parbox[c][2cm][c]{4cm}{\centering\itshape #2\\\texttt{images/#1}}}}}

\hypersetup{pdftitle={Fakerni PFE Report},pdfauthor={\myAuthor}}

%=====================================================================
\begin{document}
\pagenumbering{gobble}

%---------------------------------------------------------------------
% TITLE PAGE
%---------------------------------------------------------------------
\begin{titlepage}
\centering
{\large \mySchool\par}
\vspace{1.4cm}

\begin{minipage}{0.45\textwidth}\centering \logoph{school-logo.png}{School logo}\end{minipage}\hfill
\begin{minipage}{0.45\textwidth}\centering \logoph{company-logo.png}{Company logo}\end{minipage}

\vspace{1.6cm}
{\color{brandred}\Huge\bfseries \myReportType\par}
\vspace{0.4cm}\rule{0.7\textwidth}{0.5pt}\par
\vspace{0.8cm}
{\large In Partial Fulfillment of the Requirements for the \myDegree\par}
\vspace{0.3cm}
{\large \mySpecialization\ Specialization\par}
\vspace{0.6cm}\rule{0.7\textwidth}{0.5pt}\par

\vspace{1.0cm}
{\large Presented on : \myPresentedOn\par}
\vspace{0.4cm}
{\large Prepared by : \myAuthor\par}
\vspace{0.8cm}
{\large\itshape \guillemotleft\ \myProjectTitle\ \guillemotright\par}
\vspace{0.6cm}
{\large \myPeriod\par}

\vspace{1.4cm}
\begin{flushleft}
\hspace{1.5cm}\textbf{Professional Supervisor :} \myProSupervisor\par
\vspace{0.3cm}
\hspace{1.5cm}\textbf{Academic Supervisor :} \myAcadSupervisor\par
\end{flushleft}

\vfill
{\large ACADEMIC YEAR \myAcademicYear\par}
\end{titlepage}

%---------------------------------------------------------------------
% APPROVAL & SIGNATURES
%---------------------------------------------------------------------
\chapter*{Approval and Signatures}
\addcontentsline{toc}{chapter}{Approval and Signatures}
\vspace{1cm}
\renewcommand{\arraystretch}{3}
\begin{center}
\begin{tabular}{|p{6cm}|p{6cm}|}
\hline
\textbf{Name \& Surname} & \textbf{Date and Signature} \\ \hline
\myProSupervisor & \\ \hline
\myAcadSupervisor & \\ \hline
\end{tabular}
\end{center}
\renewcommand{\arraystretch}{1}

%---------------------------------------------------------------------
% DEDICATION
%---------------------------------------------------------------------
\chapter*{Dedication}
\addcontentsline{toc}{chapter}{Dedication}
\thispagestyle{fancy}

I dedicate this work, first and foremost, to my beloved parents, whose love,
sacrifices, and unwavering support have been the foundation of every
achievement in my life. Their patience and constant encouragement guided me
through every challenge and inspired me to pursue my goals with determination.

\medskip
To my family and my dear friends, who stood by my side throughout this journey:
thank you for your encouragement and support during both the challenging and the
rewarding moments.

\medskip
Finally, I dedicate this work to everyone who believed in me and encouraged me
throughout my academic journey. Your trust has been a driving force behind my
success.

%---------------------------------------------------------------------
% ACKNOWLEDGEMENTS
%---------------------------------------------------------------------
\chapter*{Acknowledgements}
\addcontentsline{toc}{chapter}{Acknowledgements}
\thispagestyle{fancy}

I would like to express my sincere gratitude to the entire team at \textbf{\myCompany}
for welcoming me and supporting me throughout this project. This experience was
far more than a technical assignment; it was a valuable and enriching journey
that allowed me to strengthen both my professional and personal skills.

\medskip
I extend my special thanks to my professional supervisor, \textbf{\myProSupervisor},
for the guidance, availability, and continuous support offered throughout the
project period.

\medskip
I am also deeply grateful to my academic supervisor, \textbf{\myAcadSupervisor},
for the valuable guidance, patience, and insightful recommendations that greatly
contributed to the quality and success of this work.

\medskip
My sincere appreciation also goes to all the professors and instructors of
\textbf{\mySchool} who contributed to my education and provided me with the
knowledge and skills necessary to carry out this project successfully.

%---------------------------------------------------------------------
% ABSTRACT
%---------------------------------------------------------------------
\chapter*{Abstract}
\addcontentsline{toc}{chapter}{Abstract}
\thispagestyle{fancy}

As part of my graduation project at \textbf{\myCompany}, I designed and developed a
full-stack platform named \textbf{Fakerni}, dedicated to collaborative and
real-time management of shared lists (``Fakras'') for any kind of group --- a
family, a community, an association, or an organization. The project centralizes
group membership with role-based access control, enables members to create and
share to-do and shopping lists, track budgets and spending, and coordinate in
real time across the web and mobile.

\medskip
Fakerni integrates several advanced capabilities: an AI assistant (powered by a
large language model with structured output) that turns free text and receipt
photos into structured items; a real-time collaboration layer built on
WebSockets with live presence; a centralized authorization module; budget
tracking with expense-splitting; and role-based analytics dashboards. A dynamic
multilingual interface supporting English, French, and Arabic (with
right-to-left layout) was also integrated.

\medskip
The backend was developed with \textbf{Django} and \textbf{Django REST Framework}
(with \textbf{Django Channels} for real-time), the web client with \textbf{React},
and the mobile client with \textbf{React Native (Expo)}, over a
\textbf{PostgreSQL} database, following the \textbf{Agile Scrum} methodology. This
project strengthened my skills in full-stack development, real-time systems,
role-based access control, AI integration, database design, and agile project
management.

\medskip
\noindent\textbf{Keywords:} Fakerni, Shared Lists, Collaborative Platform,
Real-Time, WebSockets, Role-Based Access Control, Artificial Intelligence,
Django, React, React Native, PostgreSQL.

%---------------------------------------------------------------------
% TABLES OF CONTENTS / FIGURES / TABLES
%---------------------------------------------------------------------
\tableofcontents
\listoffigures
\listoftables

\clearpage
\pagenumbering{arabic}

%=====================================================================
% GENERAL INTRODUCTION
%=====================================================================
\chapter*{General Introduction}
\addcontentsline{toc}{chapter}{General Introduction}

In an increasingly connected world, the coordination of everyday tasks within a
group --- deciding what to buy, who does what, and how much it costs --- still
largely relies on scattered tools: paper notes, group chats, and spreadsheets.
These tools fragment information, offer no real-time synchronization, and provide
no shared view of budgets or responsibilities. It is within this context that the
\textbf{Fakerni} project was initiated and developed at \textbf{\myCompany}.

The primary objective of the Fakerni application is to design and implement a
centralized, cross-platform system that enables:

\begin{itemize}[nosep]
  \item real-time creation and sharing of collaborative lists (``Fakras'') among
        the members of a group;
  \item structured management of groups, members, and their roles
        (owner, admin, member) with fine-grained permissions;
  \item budget tracking, spending analytics, and fair expense-splitting;
  \item AI-assisted item entry from natural language and receipt photos;
  \item role-based dashboards giving each user a tailored overview of activity.
\end{itemize}

The application consolidates all of a group's coordination data into a single
source of truth served identically to a web client and a mobile client. It
distinguishes itself through its real-time collaboration layer, its centralized
authorization model, and its practical AI features.

In accordance with modern technological standards, the application was developed
using \textbf{Django / Django REST Framework} and \textbf{Django Channels} for a
secure, high-performance, real-time backend; \textbf{React} for a responsive web
interface; \textbf{React Native (Expo)} for the mobile application; and
\textbf{PostgreSQL} as the relational database ensuring data integrity and
persistence.

The project was carried out following the \textbf{Agile Scrum} methodology,
allowing iterative and flexible development through successive sprints and
continuous adaptation to requirements. This report presents the various stages of
the project, from the initial requirements analysis and technical design through
development, testing, and the integration of user-friendly, high-performance
interfaces.

%=====================================================================
% CHAPTER 1 — GENERAL PROJECT CONTEXT
%=====================================================================
\chapter{General Project Context}

\section*{Introduction}
Before exploring the technical details, it is essential to understand the
project's environment and the business needs it addresses. This chapter
introduces the host organization, analyses the current situation and its
limitations, states the problem, presents the proposed solution, and describes
the project-management methodology adopted.

\section{Presentation of the Host Organization}
\subsection{General Presentation}
\textbf{\myCompany} is [ \textit{editable placeholder --- describe the company:
its sector, main activities, and market position} ]. The company relies on modern
technologies and qualified professionals to deliver high-quality products and
services.

Its main activities include:
\begin{itemize}[nosep]
  \item [ activity 1 ];
  \item [ activity 2 ];
  \item [ activity 3 ].
\end{itemize}

\subsection{Company Structure}
The internal structure of \myCompany\ is organized around several strategic and
operational departments that work together to ensure the efficiency of its
activities. [ \textit{Editable placeholder --- list the main departments, e.g.
Management, Development / IT, Marketing, Administration, etc.} ]

\section{Project Overview}
\subsection{Current Situation}
Currently, groups (families, roommates, associations, teams) coordinate their
shared tasks and purchases using dispersed tools such as messaging apps, paper
lists, and spreadsheets. This makes coordination complex, limits visibility over
who is responsible for what, and provides no shared, real-time view of spending.

\subsection{Analysis of the Current System}
The existing ad-hoc approach presents the following limitations:
\begin{itemize}[nosep]
  \item Lack of a centralized, shared list accessible to all members at once;
  \item No real-time synchronization --- changes are not seen instantly;
  \item No notion of roles or permissions within the group;
  \item No shared tracking of budgets, spending, or who paid what;
  \item Manual, error-prone entry of items with no assistance;
  \item No consolidated overview or analytics of group activity.
\end{itemize}

\subsection{Problem Statement}
Given the challenges identified above, there is a clear need for a single, secure,
cross-platform application that streamlines coordination between the members of a
group. The core problem statement of this project can therefore be formulated as
follows:

\begin{quote}\itshape
How can a centralized, real-time and AI-assisted platform be developed to enable
any group to create and share collaborative lists, coordinate tasks with
role-based permissions, track budgets and spending, and stay synchronized across
web and mobile devices?
\end{quote}

\subsection{Proposed Solution}
To address these challenges, we designed and developed a cross-platform
application named \textbf{Fakerni} with the following features:
\begin{itemize}[nosep]
  \item Group management with invitation codes and role-based access control
        (owner / admin / member);
  \item Shared lists (``Fakras'') with items, categories, attachments, due dates
        and recurrence;
  \item Real-time collaboration and live presence over WebSockets;
  \item Budget tracking, spending analytics, and fair expense-splitting;
  \item AI-assisted item entry (from text and receipt photos), suggestions, and a
        spending digest;
  \item Role-based interactive dashboards and push notifications;
  \item A multilingual interface (English, French, Arabic) with right-to-left
        support.
\end{itemize}

\section{Project Management Methodology}
\subsection{Different Project-Management Approaches}
Several approaches can organize a software project. The \textbf{Waterfall} model
follows a strict sequential process; it offers clear structure but limited
flexibility when requirements evolve. The \textbf{Spiral} model emphasizes
iterative risk analysis but is more complex and costly to manage. The
\textbf{Agile Scrum} framework divides the project into short iterations
(\emph{sprints}), each delivering a functional increment, and promotes
adaptability, continuous feedback, and incremental delivery.

\subsection{Chosen Methodology}
For the development of Fakerni, the \textbf{Agile Scrum} methodology was adopted.
This choice was motivated by the need to build the application incrementally while
continuously adapting to feedback and evolving requirements. The project was
organized into successive sprints, each focused on a coherent set of
functionalities. At the end of each sprint, a functional increment was delivered
and evaluated.

\subsection{SCRUM Role Distribution}
\begin{itemize}[nosep]
  \item \textbf{Product Owner:} \myProSupervisor\ --- represented the business
        requirements and validated the delivered functionalities.
  \item \textbf{Scrum Master:} \myAcadSupervisor\ --- ensured compliance with
        Scrum practices and facilitated the development process.
  \item \textbf{Development Team:} \myAuthor\ --- sole developer, responsible for
        requirements analysis, design, backend and frontend development, database
        implementation, testing, and deployment.
\end{itemize}

\section*{Conclusion}
This chapter presented the general context of the project, the host organization,
the identified challenges, and the proposed solution, together with the Agile
Scrum methodology adopted to manage the project. The next chapter details the
requirements specification and analysis.

%=====================================================================
% CHAPTER 2 — REQUIREMENTS SPECIFICATION AND ANALYSIS
%=====================================================================
\chapter{Requirements Specification and Analysis}

\section*{Introduction}
This chapter formally defines the functional and non-functional requirements of
Fakerni, identifies the system actors, presents the global use-case diagram, the
Scrum artefacts (Product Backlog and Sprint Planning), the software architecture,
and the development environment.

\section{Requirements Specification}
\subsection{Functional Requirements}
Based on the operational needs of the project, the following core functionalities
were identified and implemented:
\begin{itemize}[nosep]
  \item \textbf{Authentication and Access Control:} secure registration and login
        using email and password, JWT-based sessions with silent token refresh,
        password reset via one-time code, and centralized Role-Based Access
        Control (RBAC).
  \item \textbf{Group Management:} create a group, join via an invitation code,
        and manage members and their roles (owner, admin, member).
  \item \textbf{Shared List (Fakra) Management:} create, edit, share, archive,
        duplicate, and delete lists, with due dates and recurrence.
  \item \textbf{Item Management:} add, edit, complete/undo, delete, and assign
        items; attach photos; scan barcodes; search and filter.
  \item \textbf{Budget \& Analytics:} set a budget per list, track spending,
        receive over-budget alerts, view monthly analytics, and split expenses
        fairly among members.
  \item \textbf{AI-Assisted Features:} create items from natural-language text,
        parse items from a receipt photo, get item suggestions, run smart
        commands, and generate a spending digest.
  \item \textbf{Real-Time Collaboration:} instant synchronization of changes and
        live presence indicators across all connected clients.
  \item \textbf{Dashboards \& Notifications:} role-based overview dashboards with
        key indicators, and push notifications.
  \item \textbf{Multilingual Interface:} switch dynamically between English,
        French, and Arabic (right-to-left).
\end{itemize}

\subsection{Non-Functional Requirements}
\begin{itemize}[nosep]
  \item \textbf{Security:} all API endpoints are protected by JWT authentication;
        a centralized authorization module enforces RBAC; passwords are hashed;
        sensitive values (reset codes, device tokens) are encrypted at rest; and
        WebSocket connections are authenticated.
  \item \textbf{Performance:} the API uses query optimization (prefetching) and an
        asynchronous ASGI server for real-time traffic.
  \item \textbf{Scalability:} a stateless REST API and a Redis-backed channel
        layer allow the system to scale horizontally.
  \item \textbf{Usability:} a modern, responsive interface with dark mode and full
        internationalization, usable on both web and mobile.
  \item \textbf{Reliability:} automatic WebSocket reconnection, transparent token
        refresh, and an automated test suite executed by continuous integration.
  \item \textbf{Maintainability:} a modular architecture (separate Django apps,
        one centralized permissions module) eases evolution and bug-fixing.
  \item \textbf{Portability:} web and mobile clients share a single backend API.
\end{itemize}

\section{Global Use Case Diagram}
Figure~\ref{fig:global-uc} illustrates the global use-case diagram of the Fakerni
application, showing the interactions between the actors and the system.
\figph{uc-global.png}{Global Use Case Diagram of the Fakerni Application}
\label{fig:global-uc}

\section{Identification of Actors}
Table~\ref{tab:actors} presents the actors involved in Fakerni and their main
responsibilities.

\begin{table}[H]\centering
\caption{Actor Identification Table for the Fakerni Application}
\label{tab:actors}
\renewcommand{\arraystretch}{1.3}
\begin{tabularx}{\textwidth}{|p{2.6cm}|p{3.4cm}|X|}
\hline
\textbf{Actor} & \textbf{Description / Role} & \textbf{Main Responsibilities} \\ \hline
Group Owner & Creator and highest authority of a group &
Create the group; invite and remove members; assign roles; full control over the
group's lists; manage everything an admin can. \\ \hline
Group Admin & Manager appointed within a group &
Manage and moderate the group's lists and items; edit or delete any list or item;
help coordinate the group. \\ \hline
Member & Regular participant of a group &
Create lists; add, complete, and assign items; attach photos; view budgets,
analytics, and the shared dashboard. \\ \hline
Individual User & Any authenticated user &
Create and manage personal lists; share them with specific users; use AI features;
switch the interface language. \\ \hline
\end{tabularx}
\renewcommand{\arraystretch}{1}
\end{table}

\section{Product Backlog}
Table~\ref{tab:backlog} presents the Fakerni Product Backlog, consolidating the
functional requirements as user stories.

\begin{longtable}{|c|p{3.2cm}|p{7.2cm}|c|}
\caption{Fakerni Product Backlog}\label{tab:backlog}\\
\hline
\textbf{ID} & \textbf{Feature} & \textbf{User Story} & \textbf{Priority} \\ \hline
\endfirsthead
\hline \textbf{ID} & \textbf{Feature} & \textbf{User Story} & \textbf{Priority} \\ \hline
\endhead
1 & Authentication \& RBAC & As a user, I can register and log in securely, and my
access to features is restricted by my role. & High \\ \hline
2 & Group Management & As a user, I can create a group, invite members with a
code, and manage their roles. & High \\ \hline
3 & Shared Lists (Fakras) & As a member, I can create, share, archive, and
duplicate lists with due dates and recurrence. & High \\ \hline
4 & Item Management & As a member, I can add, complete, assign, and attach photos
to items, and scan barcodes. & High \\ \hline
5 & Budget \& Analytics & As a member, I can set budgets, track spending, split
expenses, and view analytics. & High \\ \hline
6 & AI Assistance & As a user, I can create items from text or a receipt photo and
get suggestions and a spending digest. & Medium \\ \hline
7 & Real-Time Collaboration & As a member, I can see changes and who is online in
real time. & High \\ \hline
8 & Dashboards \& Notifications & As a user, I can view a role-based overview and
receive push notifications. & Medium \\ \hline
9 & Multilingual Interface & As any user, I can switch between English, French, and
Arabic at any time. & Low \\ \hline
\end{longtable}

\section{Sprint Planning}
The project was organized into six sprints, each targeting a coherent set of
functionalities based on development priority and dependencies
(Table~\ref{tab:sprints}).

\begin{table}[H]\centering
\caption{Fakerni Sprint Planning}
\label{tab:sprints}
\renewcommand{\arraystretch}{1.3}
\begin{tabularx}{\textwidth}{|c|X|}
\hline
\textbf{Sprint} & \textbf{Features Developed} \\ \hline
Sprint 1 & Authentication, user accounts, JWT sessions, password reset, and
Role-Based Access Control (owner / admin / member). \\ \hline
Sprint 2 & Groups and shared lists (Fakras): items, categories, sharing,
attachments, barcode scanning, search and filtering. \\ \hline
Sprint 3 & Budget management, spending analytics, over-budget alerts, and expense
splitting. \\ \hline
Sprint 4 & AI-assisted features: smart add from text, receipt scanning, item
suggestions, smart commands, and spending digest. \\ \hline
Sprint 5 & Real-time collaboration: live synchronization and presence over
authenticated WebSockets. \\ \hline
Sprint 6 & Role-based dashboards and overview, notifications, and the multilingual
interface. \\ \hline
\end{tabularx}
\renewcommand{\arraystretch}{1}
\end{table}

\section{Software Architecture}
Fakerni is built on a layered client--server architecture that separates the
presentation, business-logic, and data layers, and adds a dedicated real-time
channel. Figure~\ref{fig:archi} illustrates this architecture.

\begin{itemize}[nosep]
  \item \textbf{Presentation Layer:} a \textbf{React} web client and a
        \textbf{React Native (Expo)} mobile client, both consuming the same REST
        API and connecting to the same WebSocket endpoints.
  \item \textbf{Business-Logic Layer:} a \textbf{Django / Django REST Framework}
        backend that processes requests, applies business rules, and enforces JWT
        authentication and centralized RBAC; \textbf{Django Channels} (ASGI)
        handles the real-time layer.
  \item \textbf{Data Layer:} a \textbf{PostgreSQL} relational database for
        persistent data, and \textbf{Redis} as the channel layer for real-time
        message distribution.
\end{itemize}

\figph{architecture.png}{Layered Architecture of the Fakerni Application}
\label{fig:archi}

The overall data model (entity--relationship diagram) is shown in
Figure~\ref{fig:erd}.
\figph{erd.png}{Entity--Relationship Diagram of the Fakerni Database}
\label{fig:erd}

\section{Development Environment}
\subsection{Hardware Environment}
\begin{table}[H]\centering
\caption{Hardware Environment}
\begin{tabular}{|l|l|}
\hline
\multicolumn{2}{|c|}{\textbf{Development Machine}} \\ \hline
\textbf{Processor} & [ e.g. Intel Core i5 ] \\ \hline
\textbf{RAM} & [ e.g. 16 GB ] \\ \hline
\textbf{Operating System} & [ e.g. Windows 11 ] \\ \hline
\end{tabular}
\end{table}

\subsection{Software Environment}
The main technologies and tools used to develop Fakerni are:
\begin{itemize}[nosep]
  \item \textbf{Django \& Django REST Framework} --- Python backend and REST API.
  \item \textbf{Django Channels \& Daphne} --- ASGI server and WebSockets for the
        real-time layer.
  \item \textbf{PostgreSQL} --- relational database management system.
  \item \textbf{Redis} --- channel layer for real-time message distribution.
  \item \textbf{React \& Vite} --- web front-end (single-page application).
  \item \textbf{React Native \& Expo} --- cross-platform mobile application.
  \item \textbf{Google Gemini} --- large language model used for the AI features,
        with structured (schema-constrained) output.
  \item \textbf{JSON Web Tokens (JWT)} --- stateless authentication.
  \item \textbf{Docker \& GitHub Actions} --- containerization and continuous
        integration.
  \item \textbf{Visual Studio Code \& Git / GitHub} --- code editor and version
        control.
\end{itemize}

\section{Gantt Chart}
The project timeline is structured across the six sequential sprints described
above. Figure~\ref{fig:gantt} illustrates their scheduling over the development
cycle.
\figph{gantt.png}{Fakerni Project Gantt Chart}
\label{fig:gantt}

\section*{Conclusion}
This chapter identified the functional and non-functional requirements, defined
the actors, and detailed the Product Backlog, Sprint Planning, architecture, and
development environment. The following chapters present the realization of each
sprint.

%=====================================================================
% Reusable sprint macro-sections written out explicitly per sprint
%=====================================================================

%---------------------------------------------------------------------
% CHAPTER 3 — SPRINT 1
%---------------------------------------------------------------------
\chapter{Sprint 1: Authentication, Accounts and Role-Based Access Control}

\section{Introduction}
The objective of Sprint 1 is to establish a secure and robust foundation for the
platform: user accounts, authentication, and the authorization model that governs
every later feature. This sprint implements email-based registration and login,
JWT sessions with silent refresh, password reset via a one-time code, and a
centralized Role-Based Access Control (RBAC) module.

\section{Sprint Backlog}
\begin{table}[H]\centering
\caption{Sprint 1 Backlog}
\renewcommand{\arraystretch}{1.3}
\begin{tabularx}{\textwidth}{|p{3.2cm}|X|c|c|}
\hline
\textbf{Feature} & \textbf{User Story} & \textbf{Est. (d)} & \textbf{Priority} \\ \hline
Registration & As a user, I can create an account with my email and password. & 2 & High \\ \hline
Login (JWT) & As a user, I can log in securely and stay logged in via token refresh. & 3 & High \\ \hline
Password Reset & As a user, I can reset my password using a one-time code sent to me. & 2 & Medium \\ \hline
Roles \& RBAC & As a system, access to each action is restricted by the user's role. & 3 & High \\ \hline
Profile & As a user, I can view and update my profile and device tokens. & 1 & Medium \\ \hline
\end{tabularx}
\renewcommand{\arraystretch}{1}
\end{table}

\section{Sprint 1 Analysis}
\subsection{Refinement of the Use Case ``Authenticate''}
\figph{uc-authenticate.png}{Refinement Use Case Diagram: ``Authenticate''}

\begin{table}[H]\centering
\caption{Use Case Description: ``Authenticate''}
\renewcommand{\arraystretch}{1.3}
\begin{tabularx}{\textwidth}{|p{2.8cm}|X|}
\hline
\textbf{Actor} & Any user (Owner, Admin, Member, Individual User) \\ \hline
\textbf{Preconditions} & The user has a registered account with valid credentials. \\ \hline
\textbf{Postconditions} & The user is authenticated; an access and a refresh JWT are issued and stored on the client. \\ \hline
\textbf{Main Scenario} & 1. The user enters their email and password.
2. The client sends the credentials to the API.
3. The backend verifies the hashed password and issues an access token (short-lived) and a refresh token (long-lived).
4. The client stores the tokens and attaches the access token to every request.
5. When the access token expires, the client silently obtains a new one using the refresh token. \\ \hline
\textbf{Exceptions} & Invalid credentials $\rightarrow$ ``Invalid email or password''; expired refresh token $\rightarrow$ the user is logged out and redirected to login. \\ \hline
\end{tabularx}
\renewcommand{\arraystretch}{1}
\end{table}

\subsection{Refinement of the Use Case ``Role-Based Access Control''}
Every ``who is allowed to do X'' decision is centralized in a single
authorization module. A user relates to a list through three relationships:
being its \emph{creator}, holding a \emph{group role} (owner / admin / member), or
having been granted \emph{shared access}. The backend enforces these rules on
every request; the client only uses them to hide actions the user cannot perform.

\figph{uc-rbac.png}{Refinement Use Case Diagram: ``Role-Based Access Control''}

\section{Sprint Design}
\subsection{Sequence Diagram: ``Login and Token Refresh''}
\figph{seq-login.png}{Sequence Diagram: ``Login and Token Refresh''}

\subsection{Class Diagram}
\figph{class-sprint1.png}{Class Diagram --- Sprint 1}

\subsection{Database Schema}
The main entities introduced in Sprint 1 are:
\begin{itemize}[nosep]
  \item \textbf{USER:} (\underline{id}, email, name, password (hashed), is\_active, created\_at).
  \item \textbf{DEVICE\_TOKEN:} (\underline{id}, user\_id, token (encrypted), token\_hash, created\_at).
  \item \textbf{PASSWORD\_RESET\_OTP:} (\underline{id}, user\_id, code (encrypted), expires\_at).
  \item \textbf{Roles} are carried by the \texttt{Membership} entity (owner / admin / member), introduced with groups in Sprint 2.
\end{itemize}

\section{Realization and Implementation}
The login interface authenticates the user with email and password and stores the
returned tokens. The access token is automatically attached to every request; on
expiry it is silently refreshed, so the user is never unexpectedly logged out.
Figure~\ref{fig:web-login} shows the login screen.
\figph{web-login.png}{Fakerni --- Login Interface (Web)}
\label{fig:web-login}

\section{Conclusion}
Sprint 1 delivered a secure authentication foundation and a centralized
authorization model. The next sprint builds the core domain on top of it: groups
and shared lists.

%---------------------------------------------------------------------
% CHAPTER 4 — SPRINT 2
%---------------------------------------------------------------------
\chapter{Sprint 2: Groups and Shared Lists (Fakras)}

\section{Introduction}
Building on the secure foundation of Sprint 1, Sprint 2 implements the core domain
of Fakerni: groups (with invitation codes and roles) and the shared lists
(``Fakras'') they contain --- including items, categories, attachments, barcode
scanning, sharing, archiving, duplication, and search.

\section{Sprint Backlog}
\begin{table}[H]\centering
\caption{Sprint 2 Backlog}
\renewcommand{\arraystretch}{1.3}
\begin{tabularx}{\textwidth}{|p{3.2cm}|X|c|c|}
\hline
\textbf{Feature} & \textbf{User Story} & \textbf{Est. (d)} & \textbf{Priority} \\ \hline
Groups & As a user, I can create a group and invite members with a code. & 3 & High \\ \hline
Shared Lists & As a member, I can create, share, archive, and duplicate a list. & 3 & High \\ \hline
Items & As a member, I can add, complete, assign, and edit items. & 3 & High \\ \hline
Attachments \& Barcode & As a member, I can attach a photo and scan a barcode. & 2 & Medium \\ \hline
Search \& Filter & As a member, I can search and filter items in a list. & 1 & Medium \\ \hline
\end{tabularx}
\renewcommand{\arraystretch}{1}
\end{table}

\section{Sprint 2 Analysis}
\subsection{Refinement of the Use Case ``Manage Groups and Members''}
A group has exactly one owner (enforced by a database constraint), optional
admins, and members. Membership is unique per (user, group). New members join via
a time-limited invitation code.
\figph{uc-groups.png}{Refinement Use Case Diagram: ``Manage Groups and Members''}

\subsection{Refinement of the Use Case ``Manage Shared Lists and Items''}
\begin{table}[H]\centering
\caption{Use Case Description: ``Manage Shared Lists and Items''}
\renewcommand{\arraystretch}{1.3}
\begin{tabularx}{\textwidth}{|p{2.8cm}|X|}
\hline
\textbf{Actors} & Owner, Admin, Member, Individual User \\ \hline
\textbf{Preconditions} & The user is authenticated and can view the list (creator, group member, or shared user). \\ \hline
\textbf{Postconditions} & Lists and items are created, edited, completed, or deleted, and changes are reflected for all members. \\ \hline
\textbf{Main Scenario} & 1. The user opens a list.
2. Any member can add an item or mark any item as done.
3. The item's creator, the list's creator, or a group admin can edit or delete an item.
4. The list's creator or a group owner/admin can edit, archive, share, or delete the list.
5. The backend validates every action against the centralized permissions and persists the change. \\ \hline
\textbf{Exceptions} & Insufficient permissions $\rightarrow$ HTTP 403; missing required fields $\rightarrow$ validation error. \\ \hline
\end{tabularx}
\renewcommand{\arraystretch}{1}
\end{table}
\figph{uc-lists.png}{Refinement Use Case Diagram: ``Manage Shared Lists and Items''}

\section{Sprint Design}
\subsection{Sequence Diagram: ``Add an Item to a List''}
\figph{seq-add-item.png}{Sequence Diagram: ``Add an Item to a List''}

\subsection{Class Diagram}
\figph{class-sprint2.png}{Class Diagram --- Sprint 2}

\subsection{Database Schema}
The main entities introduced in Sprint 2 are:
\begin{itemize}[nosep]
  \item \textbf{HOUSEHOLD (Group):} (\underline{id}, name, type, owner\_id, invite\_code, invite\_expires\_at, created\_at).
  \item \textbf{MEMBERSHIP:} (\underline{id}, user\_id, household\_id, role, joined\_at) --- unique per (user, household); one owner per household.
  \item \textbf{FAKRA (List):} (\underline{id}, title, description, status, due\_date, recurrence, budget, household\_id (nullable), created\_by, created\_at).
  \item \textbf{ITEM:} (\underline{id}, fakra\_id, name, quantity, unit, category, notes, estimated\_price, status, assigned\_to, done\_by, done\_at, created\_by).
  \item \textbf{ITEM\_ATTACHMENT:} (\underline{id}, item\_id, file, uploaded\_by, created\_at).
  \item \textbf{FAKRA\_ACCESS:} (\underline{id}, fakra\_id, user\_id) --- explicit sharing of a personal list.
  \item \textbf{ACTIVITY\_LOG:} (\underline{id}, fakra\_id, actor\_id, action\_type, description, created\_at).
\end{itemize}

\section{Realization and Implementation}
Figure~\ref{fig:web-fakra} shows a shared list with its items on the web client,
and Figure~\ref{fig:mobile-fakra} the same on mobile.
\figph{web-fakra.png}{Fakerni --- Shared List View (Web)}
\label{fig:web-fakra}
\figph{mobile-fakra.png}{Fakerni --- Shared List View (Mobile)}
\label{fig:mobile-fakra}

\section{Conclusion}
Sprint 2 delivered the collaborative core of Fakerni: groups, roles, and shared
lists with items and attachments. The next sprint adds budget management and
spending analytics on top of this data.

%---------------------------------------------------------------------
% CHAPTER 5 — SPRINT 3
%---------------------------------------------------------------------
\chapter{Sprint 3: Budget Management and Spending Analytics}

\section{Introduction}
Sprint 3 turns the data already collected --- items and their estimated prices ---
into financial intelligence: per-list budgets, spending tracking, over-budget
alerts, monthly analytics, and fair expense-splitting among group members.

\section{Sprint Backlog}
\begin{table}[H]\centering
\caption{Sprint 3 Backlog}
\renewcommand{\arraystretch}{1.3}
\begin{tabularx}{\textwidth}{|p{3.2cm}|X|c|c|}
\hline
\textbf{Feature} & \textbf{User Story} & \textbf{Est. (d)} & \textbf{Priority} \\ \hline
Budgets & As a member, I can set a budget for a list and track spending against it. & 2 & High \\ \hline
Over-budget Alerts & As a member, I am alerted once when spending crosses the budget. & 1 & Medium \\ \hline
Analytics & As a member, I can view monthly spending and category breakdowns. & 2 & High \\ \hline
Expense Splitting & As a group, we can compute a fair settlement of who owes whom. & 3 & High \\ \hline
\end{tabularx}
\renewcommand{\arraystretch}{1}
\end{table}

\section{Sprint 3 Analysis}
\subsection{Refinement of the Use Case ``Track Budget and Split Expenses''}
The budget cap is the explicit budget when set, otherwise the sum of item prices.
Spending is the total of completed items. An over-budget alert fires exactly once
when spending crosses the cap and resets if it drops back under. The settlement
algorithm computes each member's fair share and proposes the minimum set of
transfers to balance the group.
\figph{uc-budget.png}{Refinement Use Case Diagram: ``Track Budget and Split Expenses''}

\begin{table}[H]\centering
\caption{Use Case Description: ``Track Budget and Split Expenses''}
\renewcommand{\arraystretch}{1.3}
\begin{tabularx}{\textwidth}{|p{2.8cm}|X|}
\hline
\textbf{Actors} & Owner, Admin, Member \\ \hline
\textbf{Preconditions} & The user can view the list; items may carry estimated prices. \\ \hline
\textbf{Postconditions} & Budget indicators are updated; an alert may be sent; a settlement plan is produced. \\ \hline
\textbf{Main Scenario} & 1. A member sets a budget on the list.
2. As items are completed, spending is recomputed.
3. When spending crosses the budget, a one-time alert is sent to the members.
4. The group opens the analytics view to see spending by month and category.
5. The system computes each member's fair share and suggests the minimum transfers to settle up. \\ \hline
\textbf{Exceptions} & Negative price $\rightarrow$ validation error; no completed items $\rightarrow$ empty analytics state. \\ \hline
\end{tabularx}
\renewcommand{\arraystretch}{1}
\end{table}

\section{Sprint Design}
\subsection{Sequence Diagram: ``Complete an Item and Check the Budget''}
\figph{seq-budget.png}{Sequence Diagram: ``Complete an Item and Check the Budget''}

\subsection{Class Diagram}
\figph{class-sprint3.png}{Class Diagram --- Sprint 3}

\section{Realization and Implementation}
The list view displays a live budget bar (spent vs.\ cap) and an over-budget
warning, while the analytics view shows monthly totals and category breakdowns.
The ``settle up'' screen presents the computed transfers between members.
\figph{budget-analytics.png}{Fakerni --- Budget Bar and Spending Analytics}

\section{Conclusion}
Sprint 3 added financial intelligence to Fakerni. The next sprint introduces
AI-assisted features that reduce the effort of entering and managing items.

%---------------------------------------------------------------------
% CHAPTER 6 — SPRINT 4
%---------------------------------------------------------------------
\chapter{Sprint 4: AI-Assisted Features}

\section{Introduction}
Sprint 4 integrates a large language model to make item entry and management
effortless. All AI calls use \emph{structured output}: the model is given a strict
schema and must return valid, typed data --- so the application never has to parse
free-form text and never receives malformed results.

\section{Sprint Backlog}
\begin{table}[H]\centering
\caption{Sprint 4 Backlog}
\renewcommand{\arraystretch}{1.3}
\begin{tabularx}{\textwidth}{|p{3.2cm}|X|c|c|}
\hline
\textbf{Feature} & \textbf{User Story} & \textbf{Est. (d)} & \textbf{Priority} \\ \hline
Smart Add (text) & As a user, I can type a sentence and get structured items. & 2 & High \\ \hline
Receipt Scan & As a user, I can photograph a receipt and get its items. & 2 & High \\ \hline
Suggestions & As a user, I get relevant item suggestions for a list. & 1 & Medium \\ \hline
Smart Commands & As a user, I can issue natural-language commands on items. & 2 & Medium \\ \hline
Spending Digest & As a user, I get an AI-written summary of my spending. & 1 & Low \\ \hline
\end{tabularx}
\renewcommand{\arraystretch}{1}
\end{table}

\section{Sprint 4 Analysis}
\subsection{Refinement of the Use Case ``Assist Item Entry with AI''}
The model receives an instruction prompt and a strict output schema. For item
extraction, it returns a list of items with typed fields (name, quantity, unit,
category, notes, estimated price). For smart commands, the action is constrained
to a fixed set (done / undo / delete), so the model can only produce actions the
system can execute. AI endpoints are rate-limited to protect cost, and every call
is wrapped so a failure returns a clean error instead of crashing.
\figph{uc-ai.png}{Refinement Use Case Diagram: ``Assist Item Entry with AI''}

\section{Sprint Design}
\subsection{Sequence Diagram: ``Smart Add from Text''}
\figph{seq-ai.png}{Sequence Diagram: ``Smart Add from Text''}

\section{Realization and Implementation}
The Smart Add field turns a sentence such as ``2 kg tomatoes, milk, and bread''
into three structured items ready to save; the receipt scanner extracts items from
a photo; and the spending digest produces a short, human-readable summary.
\figph{ai-smartadd.png}{Fakerni --- AI Smart Add and Receipt Scan}

\section{Conclusion}
Sprint 4 delivered reliable, schema-constrained AI assistance. The next sprint
makes collaboration truly live with real-time synchronization and presence.

%---------------------------------------------------------------------
% CHAPTER 7 — SPRINT 5
%---------------------------------------------------------------------
\chapter{Sprint 5: Real-Time Collaboration and Presence}

\section{Introduction}
Sprint 5 makes Fakerni collaborative in real time. Using WebSockets over Django
Channels, every change made by one member appears instantly on every other
member's screen, and a presence indicator shows who is currently viewing a list.

\section{Sprint Backlog}
\begin{table}[H]\centering
\caption{Sprint 5 Backlog}
\renewcommand{\arraystretch}{1.3}
\begin{tabularx}{\textwidth}{|p{3.2cm}|X|c|c|}
\hline
\textbf{Feature} & \textbf{User Story} & \textbf{Est. (d)} & \textbf{Priority} \\ \hline
Live Sync & As a member, I see other members' changes instantly. & 3 & High \\ \hline
Live Presence & As a member, I can see who is currently viewing a list. & 2 & High \\ \hline
Secure WS Auth & As a system, WebSocket connections are authenticated with JWT. & 2 & High \\ \hline
Auto-Reconnect & As a member, my live connection recovers automatically after a drop. & 1 & Medium \\ \hline
\end{tabularx}
\renewcommand{\arraystretch}{1}
\end{table}

\section{Sprint 5 Analysis}
\subsection{Refinement of the Use Case ``Collaborate in Real Time''}
A custom middleware authenticates each WebSocket using the JWT passed in the
connection URL, and the consumer rejects any connection whose user cannot access
the list. Three channel scopes are used --- per group, per list, and per user ---
so each event reaches exactly the right audience. When a REST action changes data,
the server broadcasts an event to the relevant channel; every connected client
updates instantly. A presence registry tracks who is viewing each list and
broadcasts the (de-duplicated) viewer list on every join and leave.
\figph{uc-realtime.png}{Refinement Use Case Diagram: ``Collaborate in Real Time''}

\begin{table}[H]\centering
\caption{Use Case Description: ``Collaborate in Real Time''}
\renewcommand{\arraystretch}{1.3}
\begin{tabularx}{\textwidth}{|p{2.8cm}|X|}
\hline
\textbf{Actors} & Owner, Admin, Member, Individual User \\ \hline
\textbf{Preconditions} & The user is authenticated and opens a list they can access. \\ \hline
\textbf{Postconditions} & The client is subscribed to live updates; presence is broadcast; every change is reflected instantly. \\ \hline
\textbf{Main Scenario} & 1. The client opens a WebSocket to the list, passing its JWT.
2. The middleware authenticates the token; the consumer verifies access.
3. The user joins the list's channel and is added to the presence registry.
4. Any change (add / complete / edit / delete) is broadcast to all subscribers.
5. On leaving, the user is removed and the new viewer list is broadcast. \\ \hline
\textbf{Exceptions} & Invalid or missing token $\rightarrow$ connection refused; connection drop $\rightarrow$ automatic reconnection after a short delay. \\ \hline
\end{tabularx}
\renewcommand{\arraystretch}{1}
\end{table}

\section{Sprint Design}
\subsection{Sequence Diagram: ``Broadcast an Item Change''}
\figph{seq-realtime.png}{Sequence Diagram: ``Broadcast an Item Change''}

\section{Realization and Implementation}
When one member checks off an item, it updates on every other member's screen
within a moment, and avatars show who else is currently in the list.
\figph{realtime-presence.png}{Fakerni --- Live Updates and Presence Avatars}

\section{Conclusion}
Sprint 5 delivered authenticated, resilient real-time collaboration --- the most
technically demanding part of the project. The final sprint consolidates the
experience with role-based dashboards and the multilingual interface.

%---------------------------------------------------------------------
% CHAPTER 8 — SPRINT 6
%---------------------------------------------------------------------
\chapter{Sprint 6: Role-Based Dashboards and Multilingual Interface}

\section{Introduction}
Sprint 6 gives each user a consolidated, role-aware overview of their activity,
adds push notifications, and finalizes the dynamic multilingual interface
(English, French, and Arabic with right-to-left layout).

\section{Sprint Backlog}
\begin{table}[H]\centering
\caption{Sprint 6 Backlog}
\renewcommand{\arraystretch}{1.3}
\begin{tabularx}{\textwidth}{|p{3.2cm}|X|c|c|}
\hline
\textbf{Feature} & \textbf{User Story} & \textbf{Est. (d)} & \textbf{Priority} \\ \hline
Overview Dashboard & As a user, I see key indicators of my groups and lists at a glance. & 2 & High \\ \hline
Notifications & As a user, I receive push notifications for relevant events. & 2 & Medium \\ \hline
Multilingual & As a user, I can switch between English, French, and Arabic. & 1 & Low \\ \hline
Right-to-Left & As an Arabic user, the interface mirrors correctly (RTL). & 1 & Low \\ \hline
\end{tabularx}
\renewcommand{\arraystretch}{1}
\end{table}

\section{Sprint 6 Analysis}
\subsection{Refinement of the Use Case ``Consult Dashboard and Switch Language''}
The dashboard aggregates the user's groups, lists, pending items, and budget
status into role-appropriate indicators. The language switcher updates all labels
in real time and stores the preference locally; selecting Arabic mirrors the
layout to right-to-left.
\figph{uc-dashboard.png}{Refinement Use Case Diagram: ``Consult Dashboard and Switch Language''}

\section{Sprint Design}
\subsection{Sequence Diagram: ``Load the Overview Dashboard''}
\figph{seq-dashboard.png}{Sequence Diagram: ``Load the Overview Dashboard''}

\section{Realization and Implementation}
Figure~\ref{fig:web-dash} shows the web overview dashboard and
Figure~\ref{fig:mobile-dash} its mobile counterpart.
\figph{web-dashboard.png}{Fakerni --- Overview Dashboard (Web)}
\label{fig:web-dash}
\figph{mobile-dashboard.png}{Fakerni --- Overview Dashboard (Mobile)}
\label{fig:mobile-dash}

\section{Conclusion}
Sprint 6 completed the user experience with role-based dashboards, notifications,
and full internationalization. Together, the six sprints deliver a complete,
role-aware, real-time, AI-assisted collaborative platform.

%=====================================================================
% GENERAL CONCLUSION AND PERSPECTIVES
%=====================================================================
\chapter*{General Conclusion and Perspectives}
\addcontentsline{toc}{chapter}{General Conclusion and Perspectives}

\section*{General Conclusion}
This report was carried out as part of a project conducted at \textbf{\myCompany}.
Its objective was the design and development of \textbf{Fakerni}, a collaborative,
real-time, and AI-assisted platform enabling any group to create and share lists,
coordinate tasks with role-based permissions, track budgets, and stay synchronized
across web and mobile.

To achieve this objective, we developed a layered client--server architecture with
a single backend serving both a React web client and a React Native mobile client.
Particular effort was devoted to securing the application through JWT
authentication and a centralized Role-Based Access Control module, to making
collaboration truly live through authenticated WebSockets and presence, and to
integrating reliable AI assistance through schema-constrained model output.

Throughout the project, the \textbf{Agile Scrum} methodology ensured adaptability
and the incremental delivery of validated features at the end of each sprint. The
use of modern technologies --- Django, Django Channels, React, React Native,
PostgreSQL, and Redis --- combined with a modular architecture, made it possible to
build a robust, scalable, and maintainable solution. The implementation required
considerable research to overcome challenges such as authenticated real-time
communication, centralized multi-role authorization, and reliable AI integration.

\section*{Perspectives}
Several improvements and extensions can be envisioned for the Fakerni platform:
\begin{itemize}[nosep]
  \item \textbf{Offline Mode:} allow members to view and edit lists offline, with
        synchronization when connectivity returns.
  \item \textbf{Multi-Server Presence:} move the presence registry to Redis so
        that live presence scales across multiple backend processes.
  \item \textbf{Advanced Analytics:} add time-based trends, forecasting, and
        cross-group comparisons.
  \item \textbf{Third-Party Integrations:} connect to calendars, payment apps, and
        online grocery services.
  \item \textbf{Automated Reports:} scheduled export of activity and spending
        summaries in PDF or spreadsheet formats.
\end{itemize}

%=====================================================================
% WEBOGRAPHY
%=====================================================================
\chapter*{Webography}
\addcontentsline{toc}{chapter}{Webography}

\begin{enumerate}[label={[\arabic*]}]
  \item \textbf{Django} --- high-level Python web framework used for the backend.
        \url{https://www.djangoproject.com/}
  \item \textbf{Django REST Framework} --- toolkit for building the REST API.
        \url{https://www.django-rest-framework.org/}
  \item \textbf{Django Channels} --- WebSockets and the real-time (ASGI) layer.
        \url{https://channels.readthedocs.io/}
  \item \textbf{React} --- JavaScript library for the web user interface.
        \url{https://react.dev/}
  \item \textbf{React Native \& Expo} --- framework for the mobile application.
        \url{https://reactnative.dev/} , \url{https://expo.dev/}
  \item \textbf{PostgreSQL} --- relational database management system.
        \url{https://www.postgresql.org/}
  \item \textbf{Redis} --- in-memory data store used as the channel layer.
        \url{https://redis.io/}
  \item \textbf{JSON Web Tokens (JWT)} --- open standard for stateless
        authentication. \url{https://jwt.io/}
  \item \textbf{Google Gemini API} --- large language model used for AI features.
        \url{https://ai.google.dev/}
  \item \textbf{Scrum Guide} --- the Agile Scrum framework reference.
        \url{https://www.scrumguides.org/}
\end{enumerate}

\end{document}
